\section{Metric Spaces}

A function $d: X \times X \to \R$ is called a \textbf{metric} on $X$ if for all $p, q, r \in X$,
\begin{enumerate}[label=\textbf{(\Alph*)}]
    \item $d(p, q) \geq 0$
    \item $d(p, q) = 0 \iff p = q$
    \item $d(p, q) = d(q, p)$
    \item $d(p, r) \leq d(p, q) + d(q, r)$
\end{enumerate}
\textbf{(D)} is called the \textbf{triangle inequality}. The pair $(X, d)$ is called a \textbf{metric space}.

\bigskip

\begin{theorem} \label{an:thm:standard-metric-on-rn}
    Let $n \in \N$ be given. Define $d: \R^{n} \times \R^{n} \to \R$ as
    \[
    d(\bm{x}, \bm{y}) \coloneqq \norm*{\bm{x} - \bm{y}} = \sqrt{\sum_{k = 1}^{n} (x_{k} - y_{k})^{2}}
    \]
    for all $\bm{x}, \bm{y} \in \R^{n}$. Then $d$ is a metric on $\R^{n}$.
    In particular, $d$ is called the \textbf{standard metric} on $\R^{n}$.
\end{theorem}

\begin{proof}
    Let $\bm{x}, \bm{y}, \bm{z} \in \R^{n}$ be given.
    \begin{itemize}
        \item $d(\bm{x}, \bm{y}) \geq 0$ is trivial.
        \item By \Cref{an:def:inner-product-norm-distance},
        \begin{align*}
            d(\bm{x}, \bm{y}) = 0 &\iff \norm*{\bm{x} - \bm{y}} = 0 \\
            &\iff \forall k \in \N^{n}, x_{k} = y_{k} \\
            &\iff \bm{x} = \bm{y}
        \end{align*}
        \item $d(\bm{x}, \bm{y}) = d(\bm{y}, \bm{x})$ is trivial.
        \item By \Cref{an:thm:properties-of-inner-product-norm-distance},
        \begin{align*}
            d(\bm{x}, \bm{z}) &= \norm*{\bm{x} - \bm{z}} =
            \norm*{(\bm{x} - \bm{y}) + (\bm{y} - \bm{z})} \\
            &\leq \norm*{\bm{x} - \bm{y}} + \norm*{\bm{y} - \bm{z}}
            = d(\bm{x},\bm{y}) + d(\bm{y},\bm{z})
        \end{align*}
    \end{itemize}
    Therefore, $d$ is a metric on $\R^{n}$.
\end{proof}

\bigskip

\begin{definition} \label{an:def:k-cell}
    Let $k \in \N$ be given. A set $E \subseteq \R^{k}$ is called a \textbf{$k$-cell} if
    \[
    E = \setbuilder*{\bm{x} \in \R^{k}}{\forall j \in \N^{k}, a_{j} \leq x_{j} \leq b_{j}}
    \]
    where
    \[
    \forall j \in \N^{k}, - \infty < a_{j} \leq b_{j} < + \infty
    \]
\end{definition}

\bigskip

\begin{definition} \label{an:def:neighborhood-open-ball}
    Let $(X, d)$ be a metric space. Let $p \in X$ and $r \in \R^{+}$ be given.
    \begin{itemize}
        \item The \textbf{neighborhood} of $p$ with radius $r$ is defined as
        \[
        N_{r}(p) \coloneqq \setbuilder*{q \in X}{d(p, q) < r}
        \]
        \item The \textbf{deleted neighborhood} of $p$ with radius $r$ is defined as
        \[
        \widetilde{N}_{r}(p) \coloneqq N_{r}(p) \setminus \set*{p} = \setbuilder*{q \in X}{0 < d(p, q) < r}
        \]
    \end{itemize}
    In particular, if
    \begin{itemize}
        \item $X = \R^{n}$
        \item $d$ is the standard metric on $\R^{n}$
    \end{itemize}
    then $N_{r}(p)$ is called the \textbf{open ball} of radius $r$ centered at $p$, which is denoted by $B_{r}(p)$.
\end{definition}

\bigskip

\begin{definition} \label{an:def:limit-isolated-interior-points}
    Let $(X, d)$ be a metric space and $E \subseteq X$. Let $p \in X$ be given.
    \begin{itemize}
        \item $p$ is called an \textbf{limit point} of $E$ if
        \[
        \forall r \in \R^{+}, E \cap \widetilde{N}_{r}(p) \neq \vnot
        \]
        The set of all limit points of $E$ is denoted by $E'$.
        \item $p$ is called a \textbf{isolated point} of $E$ if
        \[
        \exists r \in \R^{+}, E \cap N_{r}(p) = \set*{p}
        \]
        \item $p$ is called an \textbf{interior point} of $E$ if
        \[
        \exists r \in \R^{+}, N_{r}(p) \subseteq E
        \]
    \end{itemize}
\end{definition}

\bigskip

\begin{definition} \label{an:def:open-closed-sets}
    Let $(X, d)$ be a metric space and $E \subseteq X$.
    \begin{itemize}
        \item $E$ is said to be \textbf{open} if every point of $E$ is an interior point of $E$.
        \item $E$ is said to be \textbf{closed} if every limit point of $E$ belongs to $E$.
    \end{itemize}
\end{definition}

\bigskip

\begin{definition} \label{an:def:bounded-perfect-dense-sets}
    Let $(X, d)$ be a metric space and $E \subseteq X$.
    \begin{itemize}
        \item $E$ is said to be \textbf{bounded} if
        \[
        \exists p \in X, \exists r \in \R^{+}, E \subseteq N_{r}(p)
        \]
        \item $E$ is said to be \textbf{perfect} if every point of $E$ is a limit point of $E$.
        \item $E$ is said to be \textbf{dense} if every point of $X$ is either a limit point of $E$ or a point of $E$.
    \end{itemize}
\end{definition}

\bigskip

\begin{theorem} \label{an:thm:neighborhood-is-open}
    Let $(X, d)$ be a metric space.
    \[
    \forall p \in X, \forall r \in \R^{+}, N_{r}(p) \text{ is open}
    \]
\end{theorem}

\begin{proof}
    Let $q \in N_{r}(p)$ be given. Then
    \[
    \begin{array}{c}
        d(p, q) < r \\
        \Downarrow \\
        \veps \coloneqq r - d(p, q) > 0
    \end{array}
    \]
    Let $s \in N_{\veps}(q)$ be given. Then
    \[
    \begin{array}{c}
        d(q, s) < \veps \\
        \Downarrow \\
        d(p, s) \leq d(p, q) + d(q, s) < d(p, q) + \veps = r \\
        \Downarrow \\
        s \in N_{r}(p) \\
        \Downarrow \\
        N_{\veps}(q) \subseteq N_{r}(p)
    \end{array}
    \]
    Since $q$ is arbitrary, $N_{r}(p)$ is open.
\end{proof}

\bigskip

\begin{theorem} \label{an:thm:limit-point-neighborhood-contains-infinite-elements}
    Let $(X, d)$ be a metric space and $E \subseteq X$.
    \[
    \forall p \in E', \forall r \in \R^{+}, \abs*{E \cap N_{r}(p)} \geq \aleph_{0}
    \]
\end{theorem}

\begin{proof}
    Assume that
    \[
    \exists p \in E', \exists r \in \R^{+}, \abs*{E \cap N_{r}(p)} < \aleph_{0}
    \]
    Define
    \[
    S \coloneqq (E \cap N_{r}(p)) \setminus \set*{p} = E \cap \widetilde{N}_{r}(p) \subseteq E \cap N_{r}(p)
    \]
    Then
    \[
    \begin{array}{c}
        p \in E' \\
        \Downarrow \\
        S \neq \vnot
    \end{array}
    \]
    By \Cref{an:thm:subset-of-finite-set-is-finite},
    \[
    \begin{array}{c}
        \abs*{E \cap N_{r}(p)} < \aleph_{0} \\
        \Downarrow \\
        \abs*{S} < \aleph_{0}
    \end{array}
    \]
    Define
    \[
    T \coloneqq \setbuilder*{d(p, q)}{q \in S} \neq \vnot
    \]
    Then
    \[
    \begin{array}{c}
        p \notin S \\
        \Downarrow \\
        0 \notin T
    \end{array}
    \]
    Define $f: S \twoheadrightarrow T$ as
    \[
    f(q) = d(p, q)
    \]
    By \Cref{an:thm:surjection-from-finite-set-implies-finite},
    \[
    \abs*{T} < \aleph_{0}
    \]
    By \Cref{an:thm:finite-ordered-set-has-extremum},
    \[
    \exists \veps \in \R^{+}, \veps = \min T
    \]
    Then
    \[
    \begin{array}{c}
        N_{\veps}(p) \subseteq N_{r}(p) \\
        \Downarrow \\
        \forall q \in S, d(p, q) < \veps \\
        \Downarrow \\
        E \cap \widetilde{N}_{\veps}(p) = \vnot \\
        \Downarrow \\
        \bot
    \end{array}
    \]
    Therefore,
    \[
    \abs*{E \cap N_{r}(p)} \geq \aleph_{0}
    \]
\end{proof}
